%% DRAFT — Baseline Methods REWRITE. Replace Section 5.3 of Final_Report.tex.

\subsection{Baseline Methods}
\label{sec:baseline_methods}

We compare against two unsupervised anomaly detectors that operate on the same edge feature vector used by graph scoring.

\textbf{Isolation Forest}~\cite{isoforest} isolates anomalies through random recursive partitions of the feature space. Anomalous points, being few and different, require fewer splits to isolate. We use 100 trees with \texttt{contamination="auto"} and a fixed random state of 42. Anomaly scores are obtained by negating scikit-learn's \texttt{score\_samples()} output.

\textbf{One-Class SVM} learns a decision boundary that encloses normal observations in a kernel-transformed space. We use the RBF kernel with $\nu = 0.1$. Features are standardized with \texttt{StandardScaler} fitted on the normal training split before being passed to the model. Anomaly scores come from negating \texttt{decision\_function()}.

Both baselines use the same feature whitelists identified by the feature audit (Section~\ref{sec:feature_selection}), so the feature sets are identical across all methods. Non-binary features are converted to percentile ranks; binary features pass through unchanged. Self-loops and user-metadata edges are excluded. Each baseline trains only on normal edges, then evaluates on held-out normal edges combined with all red-team edges, following the standard unsupervised anomaly detection protocol. Baseline scores are fed through the same threshold optimizer and pair-metric calculator used by the graph method, keeping evaluation aligned across all approaches.
